\subsection{Trajectory Dataset}

We analyzed the published three-dimensional trajectories of 19 laboratory
\textit{Chironomus riparius} swarms sampled at 100 Hz
\cite{sinhuber2019_dataset}. Each observation contains time-resolved
individual identifiers, positions, velocities, and accelerations. The original
experimental setup, imaging system, and reconstruction procedure are described
in the dataset paper.

The observations differ in duration, number of tracked individuals, track
length, and recording context. In this paper, these differences are treated as
part of the empirical domain of the result. They are not removed by pooling,
and they are not used as causal explanations unless independently verified.
The 19 observations are separate experimental recordings of the same
laboratory swarm system rather than 19 independent biological populations. We
therefore used observation identity as the replication and grouping unit
throughout the cross-observation analyses.

\begin{figure*}[t]
\centering
\includegraphics[width=\textwidth]{figures/Fig1_final.pdf}
\caption{Data and T1 measurement definition. Panel A shows a representative
raw three-dimensional snapshot from one observation. Panels B and C define the
affine subtraction and the frozen T1 measurement target used in this study.
Panel D shows the event-aligned all-tangential profile used to illustrate the
measurement pipeline. The figure defines the observable and does not represent
a fitted physical mechanism.}
\label{fig:data_t1_definition}
\end{figure*}
